\subsection{Clusters, resolution, and physical heat time}
\label{sec:cluster-resolution}

An unresolved pair need not be assigned two persistent individual labels.
It can instead be retained as a cluster with a certified total count.
The resolution here concerns what the approximation error allows one to
infer about \(F\), not an intrinsic inability to distinguish exact zeros
of a known \(G_N\). The scale depends on local separation and boundary
modulus as well as degree. In particular it is not a universal radius
depending only on \(N\).

There are two independent parameters. Increasing \(N\) improves the
approximation at a fixed physical heat time \(\lambda\); varying
\(\lambda\) changes the actual field. We now write \(G_{N,\lambda}\)
and \(F_\lambda\) explicitly. From \eqref{eq:normalized-field},
\begin{equation}
 \partial_\lambda F_\lambda
 =-\partial_z^2F_\lambda-m_2(\lambda)F_\lambda,
 \qquad m_2(\lambda)=\int u^2\,\dd\mu_\lambda(u).
 \label{eq:normalized-heat}
\end{equation}
The scalar normalization term does not move zeros. The unnormalized
\(H_\lambda\) satisfies \(\partial_\lambda H_\lambda=-\partial_z^2H_\lambda\).
Thus increasing \(\lambda\) is backward heat time; the reversed clock
\(\tau=1/2-\lambda\) has the opposite orientation.
The degree-dependent readouts are not assumed to obey that same exact
heat equation: their nodes and weights change with \(\lambda\).
The double-exponential tail permits differentiation on compact sets of
complex \((\lambda,z)\), making \(H\) jointly entire. Since
\(Z_\lambda>0\) on the real axis, the normalized field is jointly
holomorphic near each compact real time interval and spatial disk.

\begin{proposition}[Compact convergence uniform in heat time]
\label{prop:uniform-heat-compact}
For a compact real interval \(I=[\lambda_-,\lambda_+]\) and \(R<\infty\),
\begin{equation}
 \sup_{\lambda\in I,\,|z|\le R}|F_\lambda(z)-G_{N,\lambda}(z)|
 \le \frac{Z_{\lambda_+}}{Z_{\lambda_-}}
       \sum_{k=N}^\infty\frac{m_{2k}(\lambda_+)R^{2k}}{(2k)!}
 \longrightarrow0.
 \label{eq:uniform-heat-compact}
\end{equation}
\end{proposition}
\begin{proof}
Positivity gives \(Z_\lambda\ge Z_{\lambda_-}\) and
\(m_{2k}(\lambda)\le
(Z_{\lambda_+}/Z_{\lambda_-})m_{2k}(\lambda_+)\).
Apply \eqref{eq:compact-bound}; the dominating series converges by the
complete-source tail bound. All finite Jacobi coefficients, and hence
the readouts, depend continuously on real \(\lambda\), since their
moment determinants are positive and their moments are analytic.
\end{proof}

Fix a circle \(\Gamma=\{|z-a|=\rho\}\), \(a\in\R\), whose boundary
is zero-free for every \(\lambda\in I\). Suppose its interior contains
exactly two zeros of \(F_\lambda\), counted with multiplicity.
Such a circle and interval exist locally around any isolated double
zero. Set
\begin{equation}
 \eta=\min_{\lambda\in I,\,z\in\Gamma}|F_\lambda(z)|>0,
 \quad \varepsilon_N=\sup_{\lambda\in I,\,z\in\Gamma}
          |G_{N,\lambda}(z)-F_\lambda(z)|,
 \quad t_N=\varepsilon_N/\eta.
 \label{eq:cluster-boundary-budget}
\end{equation}
For \(t_N<1\), Rouch\'e gives exactly two zeros of \(G_{N,\lambda}\)
inside throughout \(I\), even when the pair cannot be separated.
The boundary minimum is a genuine hypothesis, not supplied by an upper
error bound alone. A finite-model alternative is given below.

Use contour moments to describe the cluster without choosing root labels:
\begin{equation}
 p_k(\lambda)=\frac{1}{2\pi\ii}\int_\Gamma
       (z-a)^k\frac{\partial_zF_\lambda(z)}{F_\lambda(z)}\,\dd z,
 \qquad
 c=a+\frac{p_1}{2},\qquad D=\frac{p_2}{2}-\frac{p_1^2}{4}.
 \label{eq:cluster-coordinates}
\end{equation}
Define \(p_{k,N},c_N,D_N\) by replacing \(F_\lambda\) with
\(G_{N,\lambda}\). The argument principle says \(p_0=2\), and the
residue theorem says \(p_k=(z_1-a)^k+(z_2-a)^k\). Therefore
\begin{equation}
 c=\frac{z_1+z_2}{2},\qquad
 D=\left(\frac{z_1-z_2}{2}\right)^2,\qquad
 \{z_1,z_2\}=\{c+\sqrt D,c-\sqrt D\}.
 \label{eq:pair-discriminant}
\end{equation}
The unordered pair is independent of the square-root branch. For real
\(\lambda\), conjugation symmetry makes \(c,D\) real. Positive \(D\)
means two distinct real zeros; negative \(D\) means a conjugate nonreal
pair; \(D=0\) means one real double zero. These symmetric coordinates
remain analytic in \(\lambda\) through the collision, as is immediate
from their contour definitions and joint analyticity of the field.

\begin{proposition}[Error bounds for an unresolved pair]
\label{prop:pair-resolution}
Under the preceding hypotheses and \(t_N<1\), put
\(b_N=-\log(1-t_N)\). Uniformly on \(I\),
\begin{equation}
 |p_{k,N}-p_k|\le k\rho^k b_N\quad(k\ge1),\qquad
 |c_N-c|\le\tfrac12\rho b_N,\qquad
 |D_N-D|\le 2\rho^2 b_N=:d_N.
 \label{eq:pair-resolution}
\end{equation}
In particular \(d_N\to0\). If \(D_N>d_N\), the limiting pair is
real and distinct; if \(D_N<-d_N\), it is nonreal. The band
\(|D_N|\le d_N\) is unresolved by this estimate and implies
\(|z_1-z_2|\le2\sqrt{2d_N}\).
\end{proposition}
\begin{proof}
On a neighborhood of \(\Gamma\), the ratio \(G_{N,\lambda}/F_\lambda\)
has the single-valued logarithm given by the convergent series for
\(\log(1+w)\), with \(w=(G_{N,\lambda}-F_\lambda)/F_\lambda\).
Its modulus on the circle is at most \(b_N\). Integration by parts
on the closed contour gives
\[
 p_{k,N}-p_k=-\frac{k}{2\pi\ii}\int_\Gamma
          (z-a)^{k-1}\log(G_{N,\lambda}/F_\lambda)\,\dd z.
\]
The circle length \(2\pi\rho\) proves the first bound. Both pairs
lie inside the circle, so \(|p_1|,|p_{1,N}|\le2\rho\).
Substitution into \eqref{eq:cluster-coordinates} gives the other two
bounds. The sign conclusions follow from reality of \(D,D_N\), and
the unresolved band gives \(|D|\le2d_N\).
\end{proof}

The decreasing moment-tail allowance in \eqref{eq:uniform-heat-compact}
may replace \(\varepsilon_N\), giving a decreasing resolution budget
once that allowance is below \(\eta\).

For a comparison using finite data, suppose instead that the readout
has exactly two zeros inside throughout \(I\), and
\(\gamma_N=\min_{I\times\Gamma}|G_{N,\lambda}|>\varepsilon_N\).
Rouch\'e transfers that count to \(F_\lambda\). Interchanging the
two functions in the logarithm proves the same bounds with
\(b_N=-\log(1-\varepsilon_N/\gamma_N)\). This version requires
the model boundary margin and count, but no prior knowledge of the
limiting pair's individual positions.

\begin{proposition}[Resolution near a fixed double collision]
\label{prop:double-collision-resolution}
Suppose \(F_{\lambda_*}\) has an isolated double zero at the real
point \(a\). For the local cluster coordinates just defined,
\begin{equation}
 D(\lambda_*)=0,\qquad D'(\lambda_*)=2,\qquad
 D(\lambda)=2(\lambda-\lambda_*)+O((\lambda-\lambda_*)^2).
 \label{eq:double-collision-slope}
\end{equation}
On a sufficiently small fixed interval \(I\) about \(\lambda_*\),
take \(D'\ge\kappa>0\). Every unresolved time satisfies
\begin{equation}
 |D_N(\lambda)|\le d_N
 \quad\Longrightarrow\quad
 |\lambda-\lambda_*|\le2d_N/\kappa.
 \label{eq:collision-time-resolution}
\end{equation}
For all sufficiently large \(N\), there is a time \(\lambda_N\in I\)
at which the finite readout has a double zero in the disk. Every such
time obeys \(|\lambda_N-\lambda_*|\le d_N/\kappa\).
Uniqueness of that finite-degree time is not asserted.
\end{proposition}
\begin{proof}
At \(\lambda_*\), write \(w=z-a\) and
\(F_{\lambda_*}(z)=w^2A(z)\), with \(A\) nonvanishing on the disk.
Then \(p_1=p_2=0\). Differentiate the contour formula and integrate
by parts. Equation~\eqref{eq:normalized-heat} gives
\[
 p_2'(\lambda_*)=-\frac{2}{2\pi\ii}\int_\Gamma
       w\frac{\partial_\lambda F}{F}\,\dd z
 =\frac{2}{2\pi\ii}\int_\Gamma w\frac{F_{zz}}{F}\,\dd z=4,
\]
because \(F_{zz}/F=2/w^2+4A'/(wA)+A''/A\). The normalization term
integrates to zero. Hence \(D'=p_2'/2=2\) at the event.
Analyticity proves the expansion and allows the choice of \(\kappa\).
The mean value theorem gives
\(|D(\lambda)|\ge\kappa|\lambda-\lambda_*|\), which proves
\eqref{eq:collision-time-resolution} and the stronger bound when
\(D_N=0\). Finally \(D\) has opposite signs at the two endpoints
of \(I\). Uniform convergence makes \(D_N\) have those signs too,
so continuity supplies a zero of \(D_N\).
\end{proof}

The elementary model \((z-a)^2-2(\lambda-\lambda_*)\) displays the
same orientation: as \(\lambda\) increases, a conjugate pair meets
and separates along the real axis. In reversed time the real pair
collides and leaves the axis. Its root separation is of square-root
size in elapsed time, while \(D\) is smooth. For a double zero the
spatial resolution is therefore of square-root size in the error,
whereas the time uncertainty is of linear size in the discriminant
error. This is a local approximation theorem about a fixed isolated
event. It neither asserts that a canonical collision exists nor
excludes one, and it is shared by colliding control fields with the
same local hypotheses. Higher multiplicities require more cluster
coordinates; no uniform conclusion on expanding height/time domains
is supplied here.

For background on individual real-zero dynamics in physical heat time,
Rodgers and Tao \cite[\S4, Theorem~11, Eq.~(56)]{rodgerstao} give a
principal-value interaction law in their stated range
\(\Lambda<\lambda\le0\). That simple-real-zero regime is distinct
from passing through a double event. The contour argument above needs
no individual label or division by the vanishing root separation.
